\section{Learning from outcomes: policy gradients, GAE and PPO}
\label{sec:learning}

\subsection{Reward and objective}
For a completed episode with $T$ recorded decisions, let $Y\in\{-1,0,1\}$ be
loss, tie or win from our side. The reward is terminal only:
$r_t=0$ for $t<T-1$ and $r_{T-1}=Y$. With $\gamma=0.99$,
\begin{equation}
  G_t=\sum_{j=t}^{T-1}\gamma^{j-t}r_j=\gamma^{\,T-1-t}\,Y.
\end{equation}
With $\gamma=1$, $\E[Y]=\Prob(\mathrm{win})-\Prob(\mathrm{loss})
=2\Prob(\mathrm{win})+\Prob(\mathrm{tie})-1$. When ties are absent, as in every
recorded M-C game, maximizing $\E[Y]$ is the same as maximizing win
probability. With $\gamma<1$, the same outcome is worth less after a longer
game, so training mildly prefers shorter wins and longer losses. The promotion
gate counts raw wins. The two criteria are related but not identical.

For a fixed policy, the history-conditioned value, action value and advantage
are $V^\pi(h_t)=\E_\pi[G_t\mid h_t]$,
$Q^\pi(h_t,a)=\E_\pi[G_t\mid h_t,a_t=a]$ and $A^\pi=Q^\pi-V^\pi$. The critic
approximates $V^\pi$ through compressed features. Its output is unbounded and
is not a calibrated win probability.

\subsection{The policy-gradient identity on histories}
\begin{theorem}[Finite-horizon likelihood-ratio gradient]
\label{thm:gradient}
Assume a finite trajectory space and a finite horizon $T$, padding terminated
trajectories with zero-reward absorbing steps. Assume bounded rewards, a
differentiable policy with parameter-independent support on each history, and
environment and opponent kernels that do not depend on $\theta$. Let
$J(\theta)=\E_\theta[\sum_{t<T}\gamma^tr_t]$. Then
\begin{equation}
  \nabla_\theta J(\theta)=
  \E_\theta\Big[\sum_{t=0}^{T-1}\gamma^tG_t\,
       \nabla_\theta\log\pi_\theta(a_t\mid h_t)\Big],
\label{eq:gradient}
\end{equation}
and $G_t$ may be replaced by $G_t-B(h_t)$ for any action-independent baseline
$B$ without changing the expectation.
\end{theorem}
\begin{proof}
A trajectory's probability factors into initial, transition and policy terms.
Only the policy terms depend on $\theta$, so
$\nabla\log p_\theta(\tau)=\sum_t\nabla\log\pi_\theta(a_t\mid h_t)$.
Differentiating the finite expectation gives $\E[R(\tau)\nabla\log
p_\theta(\tau)]$. A reward earned before $a_t$ is determined by $h_t$, and
$\E[\nabla\log\pi_\theta(a_t\mid h_t)\mid h_t]=\sum_a\nabla\pi_\theta(a\mid
h_t)=\nabla1=0$. Those cross terms therefore vanish, which leaves
$\sum_{j\ge t}\gamma^jr_j=\gamma^tG_t$. The same zero-mean identity removes the
baseline.
\end{proof}

\intuition{An action gains probability when the outcome that follows it is
better than expected for that situation. The theorem needs no Markov
assumption on the hashed features, because it is stated on histories. It does
need the recorded action to have actually been sampled from $\pi_\theta$.
This is why search decisions are never PPO data.}

\caveat{The implemented surrogate averages decision samples equally and omits
the $\gamma^t$ factor. It is therefore a standard heuristic related to
Equation~\eqref{eq:gradient}, not an unbiased estimator of $\nabla J$
\citep{policygradient,ppo,gae}. Opponent kernels are not fixed in self-play
either: the \code{self} opponent is a frozen snapshot of the incumbent, which
is fixed within a batch but changes after each promotion.}

\subsection{Generalized advantage estimation}
Using the collecting critic's stored values $V_t$, and $V_T=0$,
\begin{align}
  \delta_t&=r_t+\gamma V_{t+1}-V_t,\qquad
  \widehat A_t=\delta_t+\gamma\lambda\widehat A_{t+1},\qquad \widehat A_T=0,\qquad\lambda=0.95,\\
  \widehat R_t&=\widehat A_t+V_t.
\label{eq:gae}
\end{align}
Figure~\ref{fig:gae} shows how the recursion runs backward through one
episode.

\begin{figure}[htbp]
\centering
\resizebox{\linewidth}{!}{%
\begin{tikzpicture}[x=1cm,y=1cm]
  \foreach \i/\lab in {0/{preview},1/{turn 1},2/{turn 2},3/{turn 3}}{
    \node[box,text width=20mm] (s\i) at (3.6*\i,0) {$t=\i$\\\scriptsize\lab};
    \node[note] at (3.6*\i,-0.85) {$V_\i$ stored};
  }
  \node[warmbox,text width=17mm] (end) at (14.2,0) {win\\$r_3=Y=+1$};
  \foreach \i/\j in {0/1,1/2,2/3}{ \draw[arr] (s\i) -- (s\j); }
  \draw[arr] (s3) -- (end);
  \foreach \i in {0,1,2,3}{
    \node[font=\scriptsize] (d\i) at (3.6*\i,1.35) {$\delta_\i=r_\i+\gamma V_{\the\numexpr\i+1\relax}-V_\i$};
  }
  \foreach \i/\j in {3/2,2/1,1/0}{
    \draw[-{Latex[length=2mm]},thick,draw=seriesB] (3.6*\i-0.3,2.0) to[out=150,in=30]
      node[note,above]{$\times\gamma\lambda=0.9405$} (3.6*\j+0.3,2.0);
  }
  \node[font=\scriptsize,text=muted,align=left,anchor=west] at (-1.0,-1.6)
   {$\widehat A_t=\sum_{j\ge t}(\gamma\lambda)^{j-t}\delta_j$; critic target $\widehat R_t=\widehat A_t+V_t$;
    the actor uses advantages standardized once over the whole batch.};
\end{tikzpicture}}
\caption{GAE on a four-decision episode with a terminal win. Only the last
residual contains the reward. Each earlier advantage adds its own residual to
the next advantage, discounted by $\gamma\lambda$.}
\label{fig:gae}
\end{figure}

\begin{proposition}[GAE expansion and the Monte Carlo endpoint]
\label{prop:gae}
The recurrence~\eqref{eq:gae} gives
$\widehat A_t=\sum_{j=t}^{T-1}(\gamma\lambda)^{j-t}\delta_j$. If $\lambda=1$
and $V_T=0$, then $\widehat A_t=G_t-V_t$.
\end{proposition}
\begin{proof}
Unroll the recurrence backward from $\widehat A_T=0$, or argue by induction
from $t=T-1$. With $\lambda=1$, the $+\gamma V_{j+1}$ in each residual cancels
the $-V_{j+1}$ in the next one after discounting. What remains is the
discounted rewards, $-V_t$, and $\gamma^{T-t}V_T=0$.
\end{proof}

\begin{example}[Two decisions ending in a win]
With $(V_0,V_1)=(0.20,0.30)$ and $(r_0,r_1)=(0,1)$: $\delta_1=0.70$,
$\delta_0=0.99(0.30)-0.20=0.097$, $\widehat A_1=0.70$, and
$\widehat A_0=0.097+0.9405(0.70)=0.75535$. The critic targets are
$(0.95535,1.00)$; the Monte Carlo target for $t=0$ would be 0.99.
\end{example}

\subsection{Stored probabilities and importance ratios}
\label{sec:ratios}
For a step collected by revision $\theta_{\mathrm{old}}$, PPO uses
\begin{equation}
  \rho_t(\theta)=\frac{\pi_\theta(a_t\mid x_t,\A_t)}{\pi_{\theta_{\mathrm{old}}}(a_t\mid x_t,\A_t)}
  =\exp\big(\log\pi_\theta(a_t\mid x_t,\A_t)-\mathrm{storedLogProb}_t\big),
\end{equation}
with both policies tempered as in Equation~\eqref{eq:policy}.

\begin{proposition}[Conditional importance identity]
\label{prop:importance}
At a fixed stored input and candidate set, if the old policy $p$ is positive
wherever the new policy $q$ is, then $\E_{a\sim p}[(q(a)/p(a))f(a)]=\E_{a\sim
q}[f(a)]$ for every real $f$.
\end{proposition}
\begin{proof}
Write the left side as $\sum_ap(a)\,q(a)f(a)/p(a)$ and cancel $p(a)$ on its
support. The support assumption means no term with positive $q$-mass is lost.
\end{proof}

This corrects the action choice at fixed inputs only, not the whole
trajectory distribution. It is wrong when the action was really chosen by
greedy play, a human, search, or an unknown policy. Rejections add a further
subtlety. The stored likelihood is the probability of the \emph{proposal}. If
hidden restrictions reject some proposals and only accepted ones survive, the
executed action follows the proposal policy conditioned on acceptance. The
recorder therefore keeps rejected proposals separately, and accepts a step
only after its submission is confirmed.

\subsection{Collecting-likelihood validation}
\label{sec:likelihood}
Before training, the learner reloads the parent (collecting) weights and
recomputes every stored step's tempered log-probability. Any episode with a
step where the recomputed and stored values differ by more than
$\eta=10^{-4}$ is excluded.

\begin{proposition}[What the likelihood check guarantees]
\label{prop:likelihood}
Suppose every retained step satisfies
$|\log\pi_{\theta_{\mathrm{old}}}(a_t)-\mathrm{storedLogProb}_t|\le\eta$, with
the left side recomputed from the stored inputs. Then for every $\theta$ the
ratio used in training equals the exact ratio times $e^{\epsilon_t}$ with
$|\epsilon_t|\le\eta$. For $\eta=10^{-4}$, the relative ratio error is at most
$e^{\eta}-1\approx1.00005\times10^{-4}$.
\end{proposition}
\begin{proof}
Write $\epsilon_t=\log\pi_{\theta_{\mathrm{old}}}(a_t)-\mathrm{storedLogProb}_t$.
Then $\exp(\log\pi_\theta-\mathrm{storedLogProb})
=\exp(\log\pi_\theta-\log\pi_{\theta_{\mathrm{old}}})e^{\epsilon_t}$, and
$|\epsilon_t|\le\eta$ by assumption.
\end{proof}

\intuition{The check certifies more than arithmetic. A trajectory whose
features, candidate list, checkpoint or temperature changed after recording
would fail it. The six most recent learner jobs had a maximum discrepancy of
$1.9\times10^{-6}$. The check cannot certify that the recorded inputs were
themselves a faithful observation; Section~\ref{sec:game} covers that
separately.}

\subsection{The optimized PPO loss}
Advantages are standardized once over the whole batch, before any epoch:
$\widetilde A_t=(\widehat A_t-\overline A)/(s_A+10^{-8})$ when the batch has
more than one sample and $s_A>10^{-6}$. Here $s_A$ is the population
standard deviation. Each minibatch $\mathcal B$ of size $B$ also carries phase
weights $\omega_t=w_t/\overline w_{\mathcal B}$, where $w_t$ is the preview
training weight ($1\le w\le16$) for preview steps and 1 otherwise. The
minimized loss is
\begin{equation}
  \mathcal L(\theta)=\frac1B\sum_{t\in\mathcal B}\omega_t\Big[
   -\min\!\big(\rho_t\widetilde A_t,\clip(\rho_t,0.8,1.2)\widetilde A_t\big)
   +0.5\big(V_\theta(x_t)-\widehat R_t\big)^2
   -c_H\,H\big(\pi_\theta(\cdot\mid x_t,\A_t)\big)\Big].
\label{eq:loss}
\end{equation}
The entropy is that of the tempered distribution, and the value loss is not
clipped. The deployed configuration is: entropy coefficient $c_H=0.01$; preview
weight 1 (so $\omega_t\equiv1$); Adam with learning rate $3\cdot10^{-4}$;
minibatches of 32 drawn by a seeded permutation; four epochs; global
gradient-norm clip 0.5; target KL 0.03. A non-finite loss aborts before any
checkpoint is written. Each training call produces a new revision UUID and
increments the update counter once.

\begin{proposition}[What clipping bounds]
\label{prop:clip}
For a fixed advantage $A$ and positive ratio $\rho$,
\[
  \min(\rho A,\clip(\rho,1-\epsilon,1+\epsilon)A)=
  \begin{cases}A\min(\rho,1+\epsilon),&A\ge0,\\A\max(\rho,1-\epsilon),&A<0,\end{cases}
\]
and this is never larger than the unclipped term $\rho A$.
\end{proposition}
\begin{proof}
For $A\ge0$ multiplication by $A$ preserves order, so the minimum picks the
smaller of $\rho$ and $\clip(\rho)$. That is $\min(\rho,1+\epsilon)$, because
clipping from below never lowers $\rho$. For $A<0$ the order reverses. The
minimum picks the larger of the two, which is $\max(\rho,1-\epsilon)$. Taking
a minimum with $\rho A$ cannot exceed $\rho A$.
\end{proof}

\begin{example}[A clipped update]
If $\pi_{\mathrm{old}}(a)=0.20$, $\pi_\theta(a)=0.30$ and $A=0.70$, then
$\rho=1.5$. The unclipped term is 1.05; the clipped term is
$1.2(0.70)=0.84$. With $A=-0.70$ and $\rho=0.5$ the clipped term is
$0.8(-0.70)=-0.56$, not $-0.35$.
\end{example}

The proposition does not keep every ratio within $[0.8,1.2]$, does not enforce
a KL constraint, and does not lower-bound win-rate improvement
\citep{ppo,trpo}. The KL stop below is the implementation's coarse trust
region.

\subsection{The KL early stop}
After each full epoch, the learner evaluates the whole batch at the current
parameters and computes
\begin{equation}
  \widehat{\KL}=\frac1N\sum_{i=1}^N\big[(\rho_i-1)-\log\rho_i\big],
  \label{eq:k3}
\end{equation}
where $\rho_i$ is the ratio at the stored action. If $\widehat{\KL}>0.03$, the
remaining epochs are skipped. The updates from the epoch that crossed the
threshold are \emph{kept}.

\begin{proposition}[The KL estimator]
\label{prop:k3}
Fix a state, an old policy $p$ and a new policy $q$ with the same support, and
let $\rho(a)=q(a)/p(a)$. Then each term $(\rho-1)-\log\rho$ is non-negative,
and $\E_{a\sim p}[(\rho(a)-1)-\log\rho(a)]=\KL(p\Vert q)$. Consequently, if
the stored actions were sampled from the old policy, Equation~\eqref{eq:k3} is
an unbiased estimate of the average $\KL(\pi_{\mathrm{old}}\Vert\pi_\theta)$
over the batch's stored states.
\end{proposition}
\begin{proof}
For $u>0$, $u-1-\log u\ge0$, with equality only at $u=1$, because $\log u\le
u-1$. Next, $\E_p[\rho]=\sum_aq(a)=1$, so $\E_p[\rho-1]=0$, and
$\E_p[-\log\rho]=\sum_ap(a)\log(p(a)/q(a))=\KL(p\Vert q)$. Averaging over
states with their stored actions preserves unbiasedness, by linearity.
\end{proof}

\caveat{The estimator is unbiased for $\KL(\text{old}\Vert\text{new})$
(\citealp{klapprox}). The check runs only between epochs, so the returned
candidate can exceed 0.03 by up to one epoch's movement. Recent production
batches stopped at final-epoch values between 0.012 and 0.017, so none
triggered the stop. The statistic is still computed on the same samples used
for the gradient, so it is a training-set diagnostic rather than a held-out
one.}

\subsection{Demonstrations}
For explicitly labelled demonstrations $(x_t,\A_t,a_t^*)$, a separate training
call minimizes
\begin{equation}
  \mathcal L_{\mathrm{BC}}(\theta)=-\frac1B\sum_t\omega_t\log\pi_\theta(a_t^*\mid x_t,\A_t),
\end{equation}
the weighted conditional log-likelihood. There is no ratio, entropy or value
term, and no collecting-likelihood check. The shared state encoder still
changes, so the critic moves too. Imitating a weak teacher does not make the
learner an expert.

\source{\code{Model.train} (eligibility, GAE, likelihood validation,
standardization, loss, KL stop, saving) in \code{ml/model.py};
\code{RolloutBatcher} in \code{ml/batching.py}; \code{trajectory_issues} in
\code{ml/data_quality.py}; \code{LearningConfig} in \code{ml/continuous.py}.}
